% ============================================================================
% Fichier  : partie6/P06_C22_rapport.tex
% Titre    : Le Rapport de Test d'Intrusion
% Auteur   : MINKA MI NGUIDJOÏ Thierry Emmanuel
% Version  : 1.0 -- novembre 2026
% ============================================================================
% Sources : PTES (phase de restitution) ; CVSS 3.1 pour la cotation ;
% pratiques de rapport. Les 14 constats de l'Opération LOGONE : cas fictif,
% cohérents avec les chapitres précédents.
% ============================================================================

\chapter{Le Rapport de Test d'Intrusion}
\label{chap:rapport}
\chaptermark{Le rapport de test d'intrusion}

\epigraph{%
    « Un test d'intrusion ne vaut que par son rapport~;\\
    le reste n'est qu'une collection de preuves. »%
}{Principe du cours}

\niveauChapitre{Intermédiaire}{Toutes les parties
précédentes~; Chapitre~\ref{chap:methodes}.}

\begin{boxobjectifs}
\begin{itemize}
    \item Décrire la structure d'un rapport de test d'intrusion.
    \item Rédiger une synthèse pour la direction et un détail technique.
    \item Coter chaque constat et le hiérarchiser.
    \item Rédiger un constat complet~: description, preuve, impact,
          recommandation.
    \item Présenter les 14~constats de l'Opération LOGONE comme modèle.
\end{itemize}
\end{boxobjectifs}

\bigskip

% ============================================================================
\section*{Question directrice : comment transformer dix jours de test en un document utile à la direction et aux techniciens~?}
% ============================================================================

Le test de LiZbiZ est terminé~: 14~constats, des preuves, des notes. Il faut
maintenant un document que M.~ETEME lira en dix minutes et que M.~KOUNOU
suivra ligne à ligne pour corriger. Ce chapitre construit ce document.

\medskip
\textbf{Plan du chapitre.}
\begin{enumerate}
    \item La structure du rapport (section~\ref{sec:c22-structure}).
    \item Deux lecteurs, deux niveaux (section~\ref{sec:c22-lecteurs}).
    \item Coter et hiérarchiser (section~\ref{sec:c22-cotation}).
    \item L'anatomie d'un constat (section~\ref{sec:c22-constat}).
    \item Les 14 constats de LOGONE (section~\ref{sec:c22-logone}).
\end{enumerate}

\separateur

% ============================================================================
\section{La Structure du Rapport}
\label{sec:c22-structure}
% ============================================================================

\begin{tableaucours}{L{4.0cm}L{8.2cm}}{%
    \tableauHeader{Partie} &
    \tableauHeader{Contenu}
}
Synthèse pour la direction & Verdict, risques majeurs, recommandations
prioritaires, sans jargon \\
Contexte et périmètre & Objectifs, limites, dates, règles d'engagement \\
Démarche & Méthode suivie, outils, degré de couverture \\
Constats détaillés & Un par faiblesse, coté et documenté \\
Plan de correction & Actions priorisées, responsables, échéances \\
Annexes & Preuves, liste des actions horodatées, périmètre technique \\
\end{tableaucours}
\vspace{-0.8em}
\begin{center}{\small\itshape Tableau~22.1. Structure d'un rapport de test d'intrusion}\end{center}

% ============================================================================
\section{Deux Lecteurs, Deux Niveaux}
\label{sec:c22-lecteurs}
% ============================================================================

\begin{boxalert}
\textbf{Violation typique.} Un rapport uniquement technique, illisible pour
la direction, ou une synthèse vague, inutilisable par les techniciens.

\smallskip
\textbf{Règle~:} la synthèse parle risque et décision~; le détail parle
technique et correction. Les deux disent la même vérité, à deux niveaux.
\end{boxalert}

% ============================================================================
\section{Coter et Hiérarchiser}
\label{sec:c22-cotation}
% ============================================================================

Chaque constat reçoit une gravité. On peut s'appuyer sur le système commun
de notation des vulnérabilités (CVSS~3.1, Chapitre~\ref{chap:vulnerabilites})
pour les failles techniques, et sur une appréciation du risque métier pour
le reste. La gravité combine la facilité d'exploitation et l'impact.

\begin{tableaucours}{L{2.4cm}L{3.0cm}L{6.6cm}}{%
    \tableauHeader{Gravité} &
    \tableauHeader{Nombre} &
    \tableauHeader{Exigence de traitement}
}
Critique & 2 & Correction immédiate \\
Élevée & 4 & Sous quelques jours \\
Moyenne & 5 & Planifiée au trimestre \\
Faible & 3 & Au fil de l'eau \\
\end{tableaucours}
\vspace{-0.8em}
\begin{center}{\small\itshape Tableau~22.2. Répartition des 14 constats de LOGONE}\end{center}

% ============================================================================
\section{L'Anatomie d'un Constat}
\label{sec:c22-constat}
% ============================================================================

\begin{boxprocedure}
\textbf{Rédiger un constat}
\begin{enumerate}
    \item Identifiant et titre court.
    \item Description~: la faiblesse, en termes clairs.
    \item Preuve~: ce qui l'établit (requête, capture masquée, relevé).
    \item Impact~: ce qu'un attaquant en ferait.
    \item Gravité~: cotation justifiée.
    \item Recommandation~: la correction, concrète et vérifiable.
    \item Détection~: comment voir l'exploitation en production.
\end{enumerate}
\end{boxprocedure}

Tous les encadrés « Fiche de constat » des chapitres précédents suivent
exactement cette anatomie.

% ============================================================================
\section{Les 14 Constats de LOGONE}
\label{sec:c22-logone}
% ============================================================================

\begin{tableaucours}{L{1.2cm}L{6.4cm}L{2.0cm}L{2.2cm}}{%
    \tableauHeader{Id} &
    \tableauHeader{Constat} &
    \tableauHeader{Gravité} &
    \tableauHeader{Chapitre}
}
T-01 & Compte de service vulnérable au vol de ticket (facturation) &
Critique & \ref{chap:ad} \\
T-02 & Mot de passe administrateur local en clair dans un script partagé &
Critique & \ref{chap:privileges} \\
T-03 & Accès distant sans second facteur & Élevée & \ref{chap:vulnerabilites} \\
T-05 & Accès indirect aux colis d'autrui (portail) & Élevée &
\ref{chap:injections} \\
T-06 & Protocole de partage ancien et vulnérable & Élevée &
\ref{chap:exploitation} \\
T-09 & Exécution de code sur l'ancien portail & Élevée &
\ref{chap:exploitation} \\
T-04 & Identifiants par défaut sur des équipements & Moyenne &
\ref{chap:sansfil} \\
T-07 & Signature de partage non exigée et résolutions locales actives &
Moyenne & \ref{chap:reseau} \\
T-08 & Comptes d'administration utilisés sur des postes ordinaires & Moyenne
& \ref{chap:post} \\
T-10 & Service de test exposé sur l'Internet & Moyenne & \ref{chap:nuage} \\
T-11 & Taux de clic élevé à l'hameçonnage simulé & Moyenne &
\ref{chap:social} \\
T-12 & En-têtes de sécurité HTTP absents & Faible & \ref{chap:web} \\
T-13 & Bannières de version exposées & Faible & \ref{chap:balayage} \\
T-14 & Délai d'expiration de session trop long & Faible & \ref{chap:web} \\
\end{tableaucours}
\vspace{-0.8em}
\begin{center}{\small\itshape Tableau~22.3. Les 14 constats de l'Opération LOGONE}\end{center}

L'audit des mots de passe (constat A-01, Chapitre~\ref{chap:mdp}) est
rapporté séparément, comme mission complémentaire.

\begin{boxmetier}{Pentesteur et RSSI}
\textbf{Écrire pour l'action.} Un constat réussi est celui que le
destinataire sait corriger sans rappeler l'auteur. La recommandation est
concrète, vérifiable, et distingue la correction de fond du palliatif
immédiat.
\end{boxmetier}

\separateur

\begin{boxresume}
\textbf{Ce qu'il faut retenir de ce chapitre~:}
\begin{itemize}
    \item \textbf{Le rapport} est le livrable~; tout le reste le prépare.
    \item \textbf{Deux niveaux}~: synthèse pour la direction, détail pour les
          techniciens.
    \item \textbf{Chaque constat} est coté et hiérarchisé.
    \item \textbf{L'anatomie}~: description, preuve, impact, gravité,
          recommandation, détection.
    \item \textbf{Les 14 constats de LOGONE} suivent ce modèle.
\end{itemize}
\end{boxresume}

\begin{boxinfo}
\textbf{Points de contrôle de l'élève.} Je sais~: donner la structure d'un
rapport~; écrire à deux niveaux~; coter et hiérarchiser~; rédiger un constat
complet~; situer chaque constat de LOGONE.
\end{boxinfo}

\bigskip

\begin{boxexo}[title={\ding{45}\quad Exercice~22.1. Structure \niveaubase}]
Citez les six parties d'un rapport de test d'intrusion.
\end{boxexo}

\begin{boxexo}[title={\ding{45}\quad Exercice~22.2. Deux lecteurs \niveaubase}]
Qu'attend la direction du rapport~? Qu'attend le technicien~?
\end{boxexo}

\begin{boxexo}[title={\ding{45}\quad Exercice~22.3. Répartition \niveaubase}]
Combien de constats critiques, élevés, moyens et faibles compte LOGONE~?
\end{boxexo}

\begin{boxexo}[title={\ding{45}\quad Exercice~22.4. Anatomie \niveaumoyen}]
Rédigez le constat T-06 selon l'anatomie en sept points.
\end{boxexo}

\begin{boxexo}[title={\ding{45}\quad Exercice~22.5. Synthèse \niveaumoyen}]
Rédigez la synthèse pour la direction de l'Opération LOGONE, en dix lignes.
\end{boxexo}

\begin{boxexo}[title={\ding{45}\quad Exercice~22.6. Priorité \niveaumoyen}]
Justifiez pourquoi T-01 et T-02 passent avant T-12 et T-13.
\end{boxexo}

\begin{boxexo}[title={\ding{45}\quad Exercice~22.7. Mission complémentaire \niveaumoyen}]
Pourquoi l'audit des mots de passe (A-01) est-il rapporté à part des
14~constats~?
\end{boxexo}

\begin{boxexo}[title={\ding{45}\quad Exercice~22.8. Recommandation \niveaumoyen}]
Réécrivez une recommandation vague (« améliorer la sécurité ») en
recommandation concrète et vérifiable.
\end{boxexo}

\begin{boxexo}[title={\ding{45}\quad Exercice~22.9. Plan de correction \niveauexpert}]
Construisez le plan de correction de LOGONE~: ordre, responsables, échéances.
\end{boxexo}

\begin{boxexo}[title={\ding{45}\quad Exercice~22.10. Rapport complet \niveauexpert}]
Rédigez le plan complet du rapport de LOGONE, partie par partie, avec le
nombre de pages visé.
\end{boxexo}

\bigskip

\begin{boxinfo}
\textbf{Pour aller plus loin.}
\begin{itemize}
    \item PTES~\cite{ptes}~; CVSS~3.1~\cite{first2019cvss31}.
\end{itemize}
\end{boxinfo}

\begin{boxinfo}
\textbf{Transition.} Un rapport sans vérification des corrections reste
théorique~: le Chapitre~\ref{chap:retest} traite du contre-test.
\end{boxinfo}
